\documentclass[handout,10pt]{beamer}

\usetheme{CambridgeUS}
\usecolortheme[RGB={0,88,127}]{structure}
\usefonttheme{professionalfonts}

\usepackage[spanish,es-noshorthands]{babel}
\usepackage[utf8]{inputenc}
\usepackage[T1]{fontenc}

\usepackage{amsmath,amssymb,amsfonts}
\usepackage{mathpazo}
\usepackage{multicol}
\usepackage{array}
\usepackage{booktabs}
\usepackage{tikz}
\usetikzlibrary{arrows.meta}

\setbeamertemplate{navigation symbols}{}

\newtheorem{definicion}{Definición}
\newtheorem{ejemplo}{Ejemplo}
\newtheorem{ejemplos}{Ejemplos}

\title[MAT-016]{Razonamiento Lógico Matemático}
\subtitle{Unidad II: Álgebra en los Números Reales\\
	Clase 7: Conjuntos Numéricos y Operatoria}
\author[]{Profesor Luis Tulio González Alcaino\\
	Magíster en Matemática}
\institute{Universidad Santo Tomás\\
	Departamento Ciencias Básicas - Talca}
\date{Semestre I - 2026}

\begin{document}
	
	\begin{frame}
		\titlepage
	\end{frame}
	
	%------------------------------------------------
	\begin{frame}{Contenidos de la clase}
		\begin{enumerate}
			\item Conjuntos Numéricos
			\item Operatoria
		\end{enumerate}
	\end{frame}
	
	%------------------------------------------------
	\begin{frame}{Resultado de Aprendizaje}
		\begin{enumerate}
			\item Resolver operatoria en los distintos conjuntos numéricos.
		\end{enumerate}
	\end{frame}
	
	%------------------------------------------------
	\begin{frame}{Conjuntos Numéricos}
		\begin{block}{}
			Se iniciará definiendo los conjuntos de números que conforman a los números reales.
			
			\begin{enumerate}
				\item \textbf{Los Números Naturales} $\mathbb{N}$
				
				También conocidos como números para contar o enteros positivos.
				\[
				\mathbb{N}=\left\{1,2,3,\ldots\right\}
				\]
				
				\pause
				
				\item \textbf{Los Números Enteros} $\mathbb{Z}$
				
				Son los números que incluyen enteros negativos, el cero y enteros positivos:
				\[
				\mathbb{Z}=\left\{\ldots,-3,-2,-1,0,1,2,3,\ldots\right\}
				\]
			\end{enumerate}
		\end{block}
	\end{frame}
	
	%------------------------------------------------
	\begin{frame}{Conjuntos Numéricos}
		\begin{block}{}
			\begin{enumerate}
				\item[3.] \textbf{Los Números Racionales} $\mathbb{Q}$
				
				Los números racionales son los números reales que se pueden expresar como razón de dos números enteros. Se denota el conjunto de los números racionales por $\mathbb{Q}$, y se escribe
				\[
				\mathbb{Q}=
				\left\{
				\frac{a}{b}:\text{ donde } a \text{ y } b \text{ son números enteros con } b\neq 0
				\right\}
				\]
				\[
				\left\{
				\ldots,-\frac{3}{2},\ldots,-\frac{5}{4},\ldots,-1,\ldots,-\frac{1}{2},\ldots,0,\ldots,\frac{1}{2},\ldots,\frac{8}{5},\ldots
				\right\}
				\]
				
				Por otro lado, su desarrollo decimal es finito o infinito periódico o semiperiódico.
				
				\textbf{Ejemplos:}
				\[
				\frac{1}{4}=0.25
				\qquad
				\frac{5}{3}=1.666\ldots = 1.\overline{6}
				\qquad
				\frac{19}{12}=1.58333\ldots = 1.58\overline{3}
				\]
			\end{enumerate}
		\end{block}
	\end{frame}
	
	%------------------------------------------------
	\begin{frame}{Conjuntos Numéricos}
		\begin{block}{}
			\begin{enumerate}
				\item[4.] \textbf{Los Números Irracionales}
				
				Son los números que no pueden expresarse como un cociente de enteros y su desarrollo decimal es infinito no periódico.
				
				\[
				\mathbb{I}=
				\left\{
				x:\text{ donde } x \text{ no es un número racional}
				\right\}
				\]
				\[
				\left\{
				\ldots,-\pi,\ldots,-\sqrt{3},\ldots,\sqrt{2},\ldots,e,\ldots
				\right\}
				\]
				
				\textbf{Ejemplos:}
				\[
				\sqrt{2}=1.414213562\ldots
				\qquad
				\pi=3.141592654\ldots
				\qquad
				e=2.718281828\ldots
				\]
			\end{enumerate}
		\end{block}
	\end{frame}
	
	%------------------------------------------------
	\begin{frame}{Conjuntos Numéricos}
		\begin{block}{}
			\begin{enumerate}
				\item[5.] \textbf{Los Números Reales} $\mathbb{R}$
				
				La unión del conjunto de los números racionales y de los números irracionales se conoce como el conjunto de los números reales, es decir,
				\[
				\mathbb{R}=\mathbb{Q}\cup\mathbb{I}
				\]
			\end{enumerate}
			
			Una manera bastante práctica de representar el conjunto de números reales es en una recta a la que se denomina recta numérica, como la que aparece en la siguiente figura:
		\end{block}
		
		\begin{center}
			\begin{tikzpicture}[x=1.2cm,y=1cm,>=Latex]
				\draw[thick,-{Latex[length=3mm]}] (-3.6,0) -- (3.6,0);
				\foreach \x/\lab in {-3/-3,-2/-2,-1/-1,0/0,1/1,2/2,3/3}{
					\draw[thick] (\x,0.12) -- (\x,-0.12);
					\node[below] at (\x,-0.12) {\lab};
				}
			\end{tikzpicture}
		\end{center}
	\end{frame}
	
	%------------------------------------------------
	\begin{frame}{Conjuntos Numéricos}
		\begin{block}{}
			\textbf{Recta numérica}
			
			\begin{center}
				\begin{tikzpicture}[x=1.2cm,y=1cm,>=Latex]
					\draw[thick,-{Latex[length=3mm]}] (-4,0) -- (4,0);
					\foreach \x/\lab in {-3/-3,-2/-2,-1/-1,0/0,1/1,2/2,3/3}{
						\draw[thick] (\x,0.12) -- (\x,-0.12);
						\node[below] at (\x,-0.12) {\lab};
					}
				\end{tikzpicture}
			\end{center}
			
			Mediante las ideas de igualdad y desigualdad pueden compararse 2 números reales cualquiera. Se utilizan símbolos para representar estas ideas:
		\end{block}
		
		\pause
		\textbf{Ejemplos}
		\begin{itemize}
			\item Para decir que 7 es menor que 12, escribimos $7<12$.
			\pause
			\item También para decir que $-3$ es mayor que $-10$, escribimos $-3>-10$.
			\pause
			\item Podemos tener claro el significado de los símbolos menor que $<$ y mayor que $>$, si recordamos que éstos siempre apuntan hacia el número más pequeño: $-5<1$, así como $15>6$.
			\pause
			\item Hay otros dos símbolos $\leq$ y $\geq$, que también representan la idea de desigualdad. El símbolo $\leq$ significa “es menor o igual que”, por lo que $10\leq 15$ significa que 10 es menor o igual a 15.
		\end{itemize}
	\end{frame}
	
	%------------------------------------------------
	\begin{frame}{Conjuntos Numéricos}
		\begin{block}{Ejercicio}
			Complete el siguiente cuadro: Marca con ``$\surd$'' el (o los) conjunto(s) numérico(s) al(los) cual(es) pertenece cada número.
			
			\[
			\begin{array}{|c|c|c|c|c|c|}
				\hline
				\text{Número} & \mathbb{N} & \mathbb{Z} & \mathbb{Q} & \mathbb{I} & \mathbb{R} \\ \hline
				-5 &  &  &  &  &  \\ \hline
				\sqrt{3} &  &  &  &  &  \\ \hline
				\frac{5}{8} &  &  &  &  &  \\ \hline
				\frac{\pi}{2} &  &  &  &  &  \\ \hline
				1,2 &  &  &  &  &  \\ \hline
			\end{array}
			\]
		\end{block}
	\end{frame}
	
	%------------------------------------------------
	\begin{frame}{Conjuntos Numéricos}
		\begin{block}{Adición}
			Para sumar números enteros, debemos tener en cuenta que:
			
			\begin{enumerate}
				\item ``Si son del mismo signo (positivos o negativos), se suman los números (sin considerar el signo) y se conserva el signo de los números (positivo o negativo)''
				
				\textbf{Ejemplos:}
				\begin{enumerate}
					\item $2+5=7$
					\item $-4+(-7)=-11$
					\pause
				\end{enumerate}
				
				\item ``Si son de distinto signo (positivos o negativos o viceversa), se restan los números y se conserva el signo del número mayor sin considerar el signo''
				
				\textbf{Ejemplos:}
				\begin{enumerate}
					\item $8+(-3)=5$
					\item $-10+9=-1$
				\end{enumerate}
			\end{enumerate}
		\end{block}
	\end{frame}
	
	%------------------------------------------------
	\begin{frame}{Conjuntos Numéricos}
		\begin{block}{Observaciones}
			\begin{enumerate}
				\item La expresión ``$a+(-b)$'' se escribe como ``$a-b$''.
				
				\textbf{Ejemplo.} $6+(-8)=6-8=-2$
				\pause
				
				\item La expresión ``$a-(-b)$'' se escribe como ``$a+b$''.
				
				\textbf{Ejemplo.} $12-(-4)=12+4=16$
				\pause
				
				\item La expresión ``$a+b$'' es equivalente con ``$b+a$'', es decir $a+b=b+a$ (propiedad conmutativa).
				
				\textbf{Ejemplo.} $(-25)+10=10+(-25)=10-25=-15$
				\pause
				
				\item La expresión ``$a+(b+c)$'' es equivalente con ``$(a+b)+c$'', es decir
				$a+(b+c)=(a+b)+c$ (propiedad asociativa).
				
				\textbf{Ejemplo.}
				\[
				((-12)+15)+(-3)=(-12)+(15+(-3))=(-12)+(15-3)=-12+12=0
				\]
			\end{enumerate}
		\end{block}
	\end{frame}
	
	%------------------------------------------------
	\begin{frame}{Conjuntos Numéricos}
		\begin{block}{Observaciones}
			\begin{enumerate}
				\item $a+0=a$.
				
				\textbf{Ejemplo.} $7+0=7$
				\pause
				
				\item $a+(-a)=(-a)+a=0$.
				
				\textbf{Ejemplo.} $12+(-12)=0$
				\pause
				
				\item $-(-a)=a$.
				
				\textbf{Ejemplo.} $-(-9)=9$
				\pause
				
				\item $-(a+b)=(-a)+(-b)$
				\pause
				
				\item \textbf{Ejemplo.} $-(4+8)=(-4)+(-8)=-12$
			\end{enumerate}
		\end{block}
	\end{frame}
	
	%------------------------------------------------
	\begin{frame}{Conjuntos Numéricos}
		\begin{block}{}
			\textbf{Ejercicios}
			
			\begin{enumerate}
				\item Determine el valor de las siguientes sumas:
				
				\begin{tabular}{llll}
					a) $(-3)+(-8)=$\pause &  & b) $(-5+3)+2=$\pause &  \\
					c) $-10+5=$\pause &  & d) $0+(-5+2)=$\pause &  \\
					e) $(-32)+48+(-10)=$\pause &  & f) $-(-14)+(-18)=$\pause &  \\
					g) $-(3+(-7))-7-3=$\pause &  & h) $-(-8+5)+(-8)=$ & 
				\end{tabular}
				
				\item Completar los espacios en blanco en cada una de las siguientes expresiones de modo que resulte una expresión verdadera.
				
				\begin{enumerate}
					\item $3+(-3)+6=4+\underline{\hspace{2cm}}$
					\pause
					\item $(-5)+(-(-5))+(-(-(-5)))=\underline{\hspace{2cm}}$
					\pause
					\item $9=9+\underline{\hspace{2cm}}$
					\pause
					\item $-8+6+\underline{\hspace{2cm}}=-8$
				\end{enumerate}
			\end{enumerate}
		\end{block}
	\end{frame}
	
	%------------------------------------------------
	\begin{frame}{Conjuntos Numéricos}
		\begin{block}{Multiplicación}
			Para multiplicar números enteros, debemos tener en cuenta la ``ley de los signos'', que dice lo siguiente:
			\pause
			
			\begin{itemize}
				\item Al multiplicar dos números enteros de igual signo, el resultado es un número positivo $(+)$.
				\pause
				\item Al multiplicar dos números enteros de distinto signo, el resultado es un número negativo $(-)$.
				\pause
			\end{itemize}
			
			Esto se puede resumir en la siguiente tabla:
			\pause
			
			\[
			\text{\large Ley de los Signos}
			\qquad
			\begin{array}{ccc}
				+\cdot + & = & + \\
				-\cdot - & = & + \\
				+\cdot - & = & - \\
				-\cdot + & = & -
			\end{array}
			\]
		\end{block}
	\end{frame}
	
	%------------------------------------------------
	\begin{frame}{Conjuntos Numéricos}
		\begin{block}{Observaciones}
			\begin{enumerate}
				\item $a\cdot b=b\cdot a$ (conmutatividad).
				
				\textbf{Ejemplo.} $(-2)\cdot 7=7\cdot (-2)=-14$
				\pause
				
				\item $a\cdot (b\cdot c)=(a\cdot b)\cdot c$ (asociatividad).
				
				\textbf{Ejemplo.} $4\cdot(12\cdot(-3))=(4\cdot12)\cdot(-3)$
				\pause
				
				\item $a\cdot1=a$ (neutro multiplicativo).
				
				\textbf{Ejemplo.} $83\cdot1=83$
				\pause
				
				\item $a\cdot(b+c)=a\cdot b+a\cdot c$ (distributividad).
				
				\textbf{Ejemplo.} $(-4)\cdot(5+7)=(-4)\cdot5+(-4)\cdot7=-48$
				\pause
				
				\item $a\cdot0=0$.
				
				\textbf{Ejemplo.} $(-120)\cdot0=0$
				\pause
				
				\item Si $a\cdot b=0$, entonces $a=0$ o $b=0$, o bien ambos son iguales a cero.
				
				\textbf{Ejemplo.}
				\[
				(x-1)\cdot(x+5)=0 \Longrightarrow (x-1)=0 \text{ o bien } (x+5)=0
				\]
				de donde $x=1$ o bien $x=-5$.
			\end{enumerate}
		\end{block}
	\end{frame}
	
	%------------------------------------------------
	\begin{frame}{Conjuntos Numéricos}
		\begin{block}{Ejercicios}
			Multiplicar:
			
			\begin{tabular}{lll}
				a) $(-5)\cdot0=$\pause &  & f) $(-3)\cdot(-7)=$\pause \\
				b) $((-10)\cdot4)\cdot5=$\pause &  & g) $3+2\cdot5=$\pause \\
				c) $((-20)+40)\cdot(-2)=$\pause &  & h) $(-2)\cdot(-(-14)+(-5))=$\pause \\
				d) $4\cdot(-5)=$\pause &  & i) $(-7+3)\cdot(-3)=$\pause \\
				e) $20\cdot((-8)-6)=$\pause &  & j) $5\cdot4-5+6\cdot3-2+7\cdot3=$
			\end{tabular}
		\end{block}
	\end{frame}
	
	%------------------------------------------------
	\begin{frame}{Conjuntos Numéricos}
		\begin{block}{Jerarquía o Prioridad de las Operaciones}
			Cuando se evalúa una expresión que contiene más de un símbolo de operación, se procede en el orden que describiremos a continuación.
			\pause
			
			Diremos que los procedimientos u operaciones que aparecen primero son los que tienen mayor jerarquía o prioridad que aquellos que aparecen después.
			\pause
			
			\begin{enumerate}
				\item Efectuar las operaciones dentro de los paréntesis.
				\pause
				\item Realizar las multiplicaciones y divisiones.
				\pause
				\item Realizar las sumas y restas.
				\pause
				\item Cuando haya dos o más operaciones de igual jerarquía, se procede de izquierda a derecha.
				\pause
			\end{enumerate}
			
			\textbf{Ejemplos}
			\begin{enumerate}
				\item $2+[3-(5-2\cdot3)\cdot4]=\pause 9$
				\pause
				\item $(5\cdot3-2)-[4\cdot(5\cdot3)-1]=\pause -46$
				\pause
			\end{enumerate}
		\end{block}
	\end{frame}
	
	%------------------------------------------------
	\begin{frame}{Operatoria en los Números Racionales}
		\begin{block}{Adición}
			Para sumar números racionales debemos tener en cuenta que:
			\pause
			
			\begin{enumerate}
				\item Si tienen denominador común, se suman los numeradores y se conserva el denominador:
				\[
				\frac{a}{b}+\frac{c}{b}=\frac{a+c}{b},
				\qquad b\neq0
				\]
				\pause
				
				\item Si tienen distinto denominador:
				\[
				\frac{a}{b}+\frac{c}{d}=\frac{a\cdot d+b\cdot c}{b\cdot d},
				\qquad b\neq0 \text{ y } d\neq0
				\]
				\pause
				
				\item Mínimo común denominador:
				\[
				\frac{3}{5}-\frac{8}{3}+\frac{6}{12}
				=
				\frac{12\cdot3-20\cdot8+5\cdot6}{60}
				\pause
				=
				\frac{36-160+30}{60}
				=
				-\frac{94}{60}
				=
				-\frac{47}{30}
				\]
			\end{enumerate}
		\end{block}
	\end{frame}
	
	%------------------------------------------------
	\begin{frame}{Conjuntos Numéricos}
		\begin{block}{Ejemplos}
			\begin{enumerate}
				\item $\dfrac{1}{2}-\dfrac{5}{2}+\dfrac{3}{2}
				=
				\dfrac{1-5+3}{2}
				\pause
				=
				-\dfrac{1}{2}$
				\pause
				
				\item $\dfrac{1}{3}-\dfrac{2}{5}
				=
				\dfrac{1\cdot5-3\cdot2}{3\cdot5}
				=
				\dfrac{5-6}{15}
				\pause
				=
				-\dfrac{1}{15}$
				\pause
				
				\item $\dfrac{4}{5}+\dfrac{8}{3}-\dfrac{7}{18}+\dfrac{9}{2}
				=
				\dfrac{18\cdot4+30\cdot8-5\cdot7+45\cdot9}{90}
				\pause
				=
				\dfrac{682}{90}
				=
				\dfrac{341}{45}$
				\pause
				
				\item $\dfrac{1}{4}+\dfrac{7}{5}-\dfrac{3}{8}
				\pause
				=
				\dfrac{51}{40}$
			\end{enumerate}
		\end{block}
	\end{frame}
	
	%------------------------------------------------
	\begin{frame}{Conjuntos Numéricos}
		\begin{block}{Observaciones}
			\begin{enumerate}
				\item $\dfrac{-a}{b}=\dfrac{a}{-b}=-\dfrac{a}{b}$
				\pause
				\item $\dfrac{-a}{-b}=\dfrac{a}{b}$
				\pause
				\item $-\left(\dfrac{a}{b}\right)=-\dfrac{a}{b}$
				\pause
				\item $\dfrac{a}{1}=a$
				\pause
				\item $\dfrac{a}{a}=1$, con $a\neq0$
				\pause
				\item $\dfrac{0}{a}=0$, con $a\neq0$
				\pause
				\item $\dfrac{a}{0}$ no existe
			\end{enumerate}
		\end{block}
	\end{frame}
	
	%------------------------------------------------
	\begin{frame}{Conjuntos Numéricos}
		\begin{block}{Multiplicación}
			Para multiplicar dos números racionales se multiplica numerador por numerador y denominador por denominador, esto es:
			\pause
			
			\[
			\frac{a}{b}\cdot\frac{c}{d}=\frac{a\cdot c}{b\cdot d},
			\qquad b\neq0,\ d\neq0
			\]
			\pause
			
			\textbf{Ejemplos.}
			\begin{enumerate}
				\item $\dfrac{1}{5}\cdot\dfrac{-3}{8}
				=
				\dfrac{1\cdot(-3)}{5\cdot8}
				\pause
				=
				\dfrac{-3}{40}
				=
				-\dfrac{3}{40}$
				\pause
				
				\item $\dfrac{3}{-8}\cdot\dfrac{5}{7}
				=
				\dfrac{3\cdot5}{(-8)\cdot7}
				\pause
				=
				\dfrac{15}{-56}
				=
				-\dfrac{15}{56}$
				\pause
				
				\item $6\cdot\dfrac{5}{11}
				=
				\dfrac{6}{1}\cdot\dfrac{5}{11}
				=
				\dfrac{6\cdot5}{1\cdot11}
				\pause
				=
				\dfrac{30}{11}$
			\end{enumerate}
		\end{block}
	\end{frame}
	
	%------------------------------------------------
	\begin{frame}{Conjuntos Numéricos}
		\begin{block}{Propiedades}
			\begin{enumerate}
				\item Dos fracciones son iguales si, y sólo si, el producto de sus extremos es igual al producto de sus medios, es decir
				\[
				\frac{a}{b}=\frac{c}{d}\Longleftrightarrow a\cdot d=b\cdot c
				\]
				
				\textbf{Ejemplo.} $\dfrac{6}{10}=\dfrac{9}{15}$, ya que $6\cdot15=10\cdot9=90$.
				
				\item Simplificar una fracción consiste en cancelar un factor común del numerador y del denominador.
				
				\textbf{Ejemplo.}
				\[
				\dfrac{8}{12}=\dfrac{2\cdot4}{3\cdot4}=\dfrac{2}{3}
				\]
				simplificando por el factor común 4.
				
				\item Amplificar una fracción consiste en multiplicar el numerador y denominador por un mismo número.
				
				\textbf{Ejemplo.}
				\[
				\dfrac{-3}{8}=\dfrac{-3\cdot4}{8\cdot4}=\dfrac{-12}{24}
				\]
				amplificando por 4.
			\end{enumerate}
		\end{block}
	\end{frame}
	
	%------------------------------------------------
	\begin{frame}{Conjuntos Numéricos}
		\begin{block}{Ejercicios}
			Multiplique los siguientes números racionales:
			
			\bigskip
			\begin{tabular}{lll}
				a) $\dfrac{3}{9}\cdot\left(\dfrac{-4}{3}\right)=\pause -\dfrac{4}{9}$ &
				&
				d) $\dfrac{2}{7}\cdot\dfrac{1}{2}\cdot(-3)=\pause -\dfrac{3}{7}$ \\[0.25cm]
				
				b) $\left(\dfrac{(-2)-1}{3}\right)\cdot\dfrac{2}{3}=\pause -\dfrac{2}{3}$ &
				&
				e) $\dfrac{(-2)\cdot6}{-3}\cdot\dfrac{(-12)(-9)}{4}=\pause 108$ \\[0.25cm]
				
				c) $\left(\dfrac{-6}{3}\right)\cdot\left(\dfrac{-4}{7}\right)\cdot\dfrac{1}{9}
				=\pause \dfrac{8}{63}$ &
				&
				f) $\left(\dfrac{-7}{8}\right)\cdot\left(\dfrac{-4}{-3}\right)\cdot\left(\dfrac{2}{14}\right)\cdot(-4)
				=\pause \dfrac{2}{3}$
			\end{tabular}
		\end{block}
	\end{frame}
	
	%------------------------------------------------
	\begin{frame}{Conjuntos Numéricos}
		\begin{block}{División}
			Para dividir dos números racionales se multiplica la primera fracción por la segunda invertida.
			
			\[
			\frac{a}{b}\div\frac{c}{d}=\frac{a}{b}\cdot\frac{d}{c}
			\]
			
			\textbf{Ejemplos.}
			\begin{enumerate}
				\item $\dfrac{3}{5}\div\dfrac{2}{9}
				=
				\dfrac{3}{5}\cdot\dfrac{9}{2}
				\pause
				=
				\dfrac{27}{10}$
				\pause
				
				\item $\dfrac{1}{2}\div\left(\dfrac{-2}{3}\right)
				=
				\dfrac{1}{2}\cdot\left(\dfrac{3}{-2}\right)
				\pause
				=
				-\dfrac{3}{4}$
				\pause
				
				\item $7\div\dfrac{14}{5}
				=
				\dfrac{7}{1}\cdot\dfrac{5}{14}
				\pause
				=
				\dfrac{5}{2}$
			\end{enumerate}
		\end{block}
	\end{frame}
	
	%------------------------------------------------
	\begin{frame}{Conjuntos Numéricos}
		\begin{block}{Ejercicios}
			Divida las siguientes expresiones, simplificando lo más posible
			
			\begin{tabular}{lll}
				a) $\dfrac{-8}{5}\div\dfrac{4}{7}=$\pause & &
				d) $\dfrac{7}{9}\div\dfrac{-4}{12}=$\pause \\[0.2cm]
				
				b) $\left(\dfrac{-1}{3}\right)\div\left(\dfrac{2}{-3}\right)=$\pause & &
				e) $\left(\dfrac{(-2)\cdot6}{-3}\right)\div\left(\dfrac{(-15)(-2)}{4}\right)=$\pause \\[0.2cm]
				
				c) $\left(\dfrac{-3}{4}\right)\div6=$\pause & &
				f) $\left(\left(\dfrac{-5}{8}\right)\cdot\left(\dfrac{-4}{-5}\right)\right)\div\left(\dfrac{12}{6}\right)=$
			\end{tabular}
		\end{block}
	\end{frame}
	
	%------------------------------------------------
	\begin{frame}{Ejercicios}
		\begin{block}{Ejemplos}
			Calcular las siguientes expresiones:
			
			\begin{multicols}{2}
				\begin{enumerate}
					\item $(-5+(-7))=$
					\item $(3+(-4)-8+12)=$
					\item $\dfrac{12(-40)(-12)}{-5(-3)-3(-3)}=$
					\item $\dfrac{(-3)(8)(-2)}{(-4)(-6)-2(-12)}=$
					\item $\dfrac{-2}{9}+\dfrac{4}{7}-\dfrac{1}{2}=$
					\item $\dfrac{\frac{2}{3}+\frac{1}{2}}{\frac{2}{2}}=$
					\item $\dfrac{\frac{1}{8}+\dfrac{1-\dfrac{1}{2}}{2}}{\frac{1}{12}}=$
					\item $4+\dfrac{3}{5}-\dfrac{7}{2}-\dfrac{4}{12}-\dfrac{1}{6}=$
					\item $\dfrac{(-5)\left(\frac{4}{7}\right)}{(-2)\left(\frac{-3}{5}\right)}=$
				\end{enumerate}
			\end{multicols}
		\end{block}
	\end{frame}
	
\end{document}